\documentclass[12pt,a4paper]{article}
\usepackage[utf8]{inputenc}
\usepackage[T1]{fontenc}
\usepackage{amsmath,amssymb,amsfonts}
\usepackage{geometry}
\usepackage{listings}
\usepackage{xcolor}
\usepackage{hyperref}
\usepackage{booktabs}
\usepackage{array}
\usepackage{tikz}
\usetikzlibrary{arrows.meta,shapes,positioning,calc}

\geometry{margin=2.5cm}
\lstset{
  basicstyle=\ttfamily\footnotesize,
  keywordstyle=\color{blue!80}\bfseries,
  commentstyle=\color{gray}\itshape,
  stringstyle=\color{red!70},
  breaklines=true,
  frame=single,
  backgroundcolor=\color{gray!5},
  numbers=left,
  numberstyle=\tiny\color{gray},
  captionpos=b
}

\title{\textbf{Report 3: Fixed-Wing Position Control}\\[0.4em]
\large \textbf{Part~I:} Complete Original System\\
\textbf{Part~II:} Every Soaring-Mode Modification and Its Impact\\[0.5em]
\normalsize PX4 Autosoaring Implementation --- Defence Technical Report}
\author{Radhouene --- Fixed-Wing Autosoaring Project}
\date{April 2026}

\begin{document}
\maketitle
\tableofcontents
\newpage

%=============================================================================
\section*{List of Symbols}
\addcontentsline{toc}{section}{List of Symbols}
%=============================================================================

A superscript $(\cdot)^*$ denotes a \emph{setpoint} (commanded value).
A dot $\dot{(\cdot)}$ denotes the time derivative.
Bold symbols ($\mathbf{v}$) denote 2-D or 3-D vectors.
A hat $\hat{(\cdot)}$ (in guidance context) denotes a unit vector.

\subsection*{Position and Velocity}
\begin{center}\begin{tabular}{lll}\toprule
Symbol & Meaning & Unit \\\midrule
$\mathbf{p} = (x, y)$ & Aircraft ground position in local NED North-East plane & m \\
$\mathbf{v}_g = (v_N, v_E)$ & Aircraft ground velocity (North, East) & m/s \\
$\mathbf{w} = (w_N, w_E)$ & Wind velocity vector (North, East) & m/s \\
$\mathbf{v}_a = \mathbf{v}_g - \mathbf{w}$ & Aircraft airspeed vector (ground velocity minus wind) & m/s \\
$V_a = \|\mathbf{v}_a\|$ & Airspeed magnitude & m/s \\
$v_{g,x}$ & Body-frame forward ground speed (from estimator) & m/s \\
$V_{g,\min}$ & Minimum forward ground speed (\texttt{FW\_GND\_SPD\_MIN}) & m/s \\
$h$ & Altitude AMSL, computed as $h = -z_{\text{NED}} + h_{\text{ref}}$ & m \\
$h_{\text{ref}}$ & EKF reference altitude (\texttt{local\_pos.ref\_alt}) & m \\
\bottomrule\end{tabular}\end{center}

\subsection*{Attitude Angles}
\begin{center}\begin{tabular}{lll}\toprule
Symbol & Meaning & Unit \\\midrule
$\phi$ & Bank (roll) angle, positive right-wing-down & rad \\
$\phi^*$ & Roll setpoint & rad \\
$\phi_{\max}$ & Maximum roll limit (\texttt{FW\_R\_LIM}) & rad \\
$\theta$ & Pitch angle, positive nose-up & rad \\
$\theta^*$ & Pitch setpoint (from TECS) & rad \\
$\psi$ & Yaw (heading) angle, measured clockwise from North & rad \\
$\mathbf{q}$ & Unit quaternion representing current attitude & --- \\
$\mathbf{q}^*$ & Attitude setpoint quaternion published to inner-loop & --- \\
$R$ & Direction Cosine Matrix (DCM) from body to NED & --- \\
\bottomrule\end{tabular}\end{center}

\subsection*{Airspeed}
\begin{center}\begin{tabular}{lll}\toprule
Symbol & Meaning & Unit \\\midrule
$V_{\text{EAS}}$ & Equivalent airspeed: corrected for air density & m/s \\
$V_{\text{TAS}}$ & True airspeed: actual speed relative to air mass & m/s \\
$k = k_{\text{EAS2TAS}} = \sqrt{\rho_0/\rho}$ & EAS-to-TAS conversion factor, $\geq 1$ & --- \\
$V^*_{\text{EAS}}$ & EAS setpoint after \texttt{adapt\_airspeed\_setpoint()} & m/s \\
$V_{\text{trim}}$ & Trim (cruise) EAS from performance model & m/s \\
$V_{\text{EAS,min}}$, $V_{\text{EAS,max}}$ & EAS flight envelope bounds from performance model & m/s \\
$V_{\text{TAS,min}}$, $V_{\text{TAS,max}}$ & Corresponding TAS bounds $= k\cdot V_{\text{EAS,min/max}}$ & m/s \\
$k_w$ & Wind-based airspeed scaling coefficient (\texttt{FW\_WIND\_ARSP\_SC}) & --- \\
\bottomrule\end{tabular}\end{center}

\subsection*{NPFG Guidance}
\begin{center}\begin{tabular}{lll}\toprule
Symbol & Meaning & Unit \\\midrule
$\hat{\mathbf{t}}$ & Unit tangent vector along the reference path & --- \\
$\hat{\mathbf{n}} = \hat{\mathbf{t}}^{\perp}$ & Unit normal to the path (points left of direction of travel) & --- \\
$e_y$ & Signed cross-track error: positive = aircraft is right of path & m \\
$\eta$ & Track-error angle between airspeed vector and desired heading & rad \\
$L$ & NPFG look-ahead distance (scales with airspeed and period) & m \\
$T_{\text{NPFG}}$ & NPFG guidance period parameter (\texttt{NPFG\_PERIOD}) & s \\
$\zeta$ & NPFG damping ratio (\texttt{NPFG\_DAMPING}) & --- \\
$a_{\text{lat}}$ & Lateral (horizontal) acceleration command from NPFG & m/s$^2$ \\
$R$ & Loiter circle radius (positive = clockwise) & m \\
$f_{\text{run}} \in [0,1]$ & NPFG can-run confidence factor (1 = full confidence) & --- \\
\bottomrule\end{tabular}\end{center}

\subsection*{TECS Interface (as seen from this module)}
\begin{center}\begin{tabular}{lll}\toprule
Symbol & Meaning & Unit \\\midrule
$h^*$ & Altitude setpoint passed to TECS & m \\
$\theta^*_{\text{TECS}}$ & Pitch setpoint returned by \texttt{get\_tecs\_pitch()} & rad \\
$\delta^*_{t,\text{TECS}}$ & Throttle returned by \texttt{get\_tecs\_thrust()} & --- \\
$n_z = 1/\cos\phi$ & Normal load factor injected into TECS for stall floor & --- \\
\bottomrule\end{tabular}\end{center}

\subsection*{adapt\_airspeed\_setpoint Notation}
\begin{center}\begin{tabular}{lll}\toprule
Symbol & Meaning & Unit \\\midrule
$V^*_{\text{EAS,in}}$ & Input EAS setpoint (from waypoint or mission) & m/s \\
$V^*_{\text{EAS,out}}$ & Slew-rate limited EAS setpoint passed to TECS & m/s \\
$V_{\text{EAS,min,wind}}$ & Wind-adjusted minimum EAS $= V_{\text{EAS,min}} + k_w\|\mathbf{w}\|$ & m/s \\
$r_V$ & Airspeed setpoint slew rate (\texttt{ASPD\_SP\_SLEW\_RATE}, default 1\,m/s$^2$) & m/s$^2$ \\
\bottomrule\end{tabular}\end{center}

\subsection*{Manual Control}
\begin{center}\begin{tabular}{lll}\toprule
Symbol & Meaning & Unit \\\midrule
$\delta_{\text{pitch}}$ & Pitch stick input, $\in[-1, +1]$ & --- \\
$\delta_{\text{throttle}}$ & Throttle stick input, $\in[-1, +1]$ & --- \\
$\delta_{\text{roll}}$ & Roll stick input, $\in[-1, +1]$ & --- \\
$\delta_{\text{yaw}}$ & Yaw stick input, $\in[-1, +1]$ & --- \\
$\dot{h}^*_{\text{stick}}$ & Height rate setpoint from pitch stick & m/s \\
$V^*_{\text{stick}}$ & Airspeed setpoint from throttle stick & m/s \\
$\psi_{\text{hold}}$ & Latched heading in heading-hold mode & rad \\
$\mathbf{p}_{\text{hold}}$ & Latched NE position in heading-hold mode & m \\
\bottomrule\end{tabular}\end{center}

\subsection*{GPS Failsafe}
\begin{center}\begin{tabular}{lll}\toprule
Symbol & Meaning & Unit \\\midrule
$T_{\text{gpsf}}$ & Fixed-bank loiter duration before descend (\texttt{NAV\_GPSF\_LT}, default 30\,s) & s \\
$\phi_{\text{gpsf}}$ & Fixed open-loop bank angle (\texttt{NAV\_GPSF\_R}, default 15$^\circ$) & deg \\
$\dot{h}_{\text{descent}}$ & Hard-coded descent rate in GPS failsafe = $-0.5$\,m/s & m/s \\
\bottomrule\end{tabular}\end{center}

\subsection*{Control Mode Flags (set by set\_control\_mode\_current)}
\begin{center}\begin{tabular}{ll}\toprule
Symbol / Name & Condition \\\midrule
\texttt{FW\_POSCTRL\_MODE\_AUTO} & Normal auto: position + position setpoint valid \\
\texttt{FW\_POSCTRL\_MODE\_AUTO\_TAKEOFF} & Takeoff setpoint type \\
\texttt{FW\_POSCTRL\_MODE\_AUTO\_LANDING\_STRAIGHT} & Land type + valid previous WP \\
\texttt{FW\_POSCTRL\_MODE\_AUTO\_LANDING\_CIRCULAR} & Land type, no previous WP \\
\texttt{FW\_POSCTRL\_MODE\_AUTO\_PATH} & Offboard with velocity setpoint \\
\texttt{FW\_POSCTRL\_MODE\_AUTO\_ALTITUDE} & GPS failsafe phase 1 (fixed-bank loiter) \\
\texttt{FW\_POSCTRL\_MODE\_AUTO\_CLIMBRATE} & GPS failsafe phase 2 (descend) \\
\texttt{FW\_POSCTRL\_MODE\_MANUAL\_ALTITUDE} & Manual + altitude hold \\
\texttt{FW\_POSCTRL\_MODE\_MANUAL\_POSITION} & Manual + heading hold \\
\bottomrule\end{tabular}\end{center}

\subsection*{Soaring-Mode Symbols (Part~II)}
\begin{center}\begin{tabular}{lll}\toprule
Symbol & Meaning & Unit \\\midrule
$\phi_{\text{cmd}}$ & Commanded bank angle for thermal (\texttt{bank\_angle\_cmd}) & deg \\
$|\phi_{\text{cmd}}|$ & Magnitude of commanded bank angle & deg \\
$R_{\min}$ & Minimum safe loiter radius at speed $V$ and bank $\phi$ & m \\
$R_{\text{eff}}$ & Effective loiter radius: $\max(R_{\text{companion}}, R_{\min})$ & m \\
$R_c$ & Companion-provided loiter radius (\texttt{loiter\_radius\_m}) & m \\
$T_{\text{ramp}}$ & Throttle ramp duration (\texttt{FW\_GLIDE\_RAMP\_T}, default 3\,s) & s \\
$t_s$ & Entry ramp start timestamp (\texttt{\_soaring\_entry\_ramp\_start\_us}) & $\mu$s \\
$t_e$ & Exit ramp start timestamp (\texttt{\_soaring\_ramp\_start\_us}) & $\mu$s \\
$\delta_{t,\text{pre}}$ & TECS throttle just before glide entry (\texttt{\_soaring\_pre\_glide\_thrust}) & --- \\
$s_{\text{exit}}(t)$ & Exit ramp scale factor $= \min((t-t_e)/T_{\text{ramp}},\,1)$ & --- \\
$h_{\min}$, $h_{\max}$ & Soaring altitude floor/ceiling (\texttt{FW\_ALT\_MIN}, \texttt{FW\_ALT\_MAX}) & m \\
$h_{\text{hyst}}$ & Altitude ceiling hysteresis band (\texttt{FW\_ALT\_HYST}) & m \\
$V^*_{\text{EAS,glide}}$ & Glide EAS setpoint injected into TECS for soaring modes & m/s \\
$a, b$ & Glide polar coefficients: sink rate $w_s = aV^2 + b$ & --- \\
$V^*_{\text{polar}} = \sqrt{b/a}$ & Best-glide EAS from polar (Mode 2) & m/s \\
$\mathbf{q}^* = \text{Euler}(\phi^*,\theta^*,\psi)$ & Mode~4 attitude setpoint quaternion & --- \\
\bottomrule\end{tabular}\end{center}

\subsection*{Mathematical Operators}
\begin{center}\begin{tabular}{lll}\toprule
Operator & Meaning & Notes \\\midrule
$\text{clamp}(x,a,b)$ & Saturate: $\max(a,\min(b,x))$ & Also \texttt{constrain()} in PX4 \\
$\text{lerp}(a,b,t)$ & Linear interpolation: $(1-t)a+tb$ & $t\in[0,1]$ \\
$\text{Euler}(\phi,\theta,\psi)$ & ZYX Euler angles to unit quaternion & Roll, pitch, yaw order \\
$\text{atan2}(y,x)$ & Four-quadrant arctangent & Returns angle in $(-\pi,\pi]$ \\
$\text{sign}(x)$ & Signum: $+1$ if $x>0$, else $-1$ & \\
$\|\mathbf{v}\|$ & Euclidean norm of vector $\mathbf{v}$ & \\
$\varepsilon$ & Small positive guard value against division by zero & \\
$\max(a,b)$, $\min(a,b)$ & Standard maximum / minimum & \\
$|x|$, \texttt{fabsf}$(x)$ & Absolute value & \\
\bottomrule\end{tabular}\end{center}

\newpage
% ============================================================
%  PART I  — ORIGINAL SYSTEM (every function)
% ============================================================
\part{Original Fixed-Wing Position Control System}

%=============================================================================
\section{Module Identity and Scheduling}
%=============================================================================

File: \texttt{src/modules/fw\_pos\_control/FixedwingPositionControl.cpp}

\textbf{Task trigger:} The module is scheduled by the
\texttt{local\_position} uORB topic; \texttt{Run()} is called
every time a new position estimate is published (nominally 50\,Hz,
$\Delta t = 20\,\text{ms}$).

\textbf{Interval guard:}
\begin{equation}
  \Delta t_{\text{eff}} = \text{clamp}\!\left(
    \frac{t_k - t_{k-1}}{10^6},\;
    \Delta t_{\min},\;\Delta t_{\max}
  \right)
  \quad [0.01\,\text{s},\;0.1\,\text{s}]
\end{equation}

The clamp prevents integration blow-up during long-pauses (start-up,
debugger break) and avoids division by zero.

\textbf{Output topics written each cycle (if not idle):}
\begin{itemize}
  \item \texttt{vehicle\_attitude\_setpoint} ($\mathbf{q}^*, \delta_t$)
  \item \texttt{tecs\_status} (debug / logging)
  \item \texttt{position\_controller\_status}
  \item \texttt{autosoaring\_status} (FMU $\rightarrow$ companion: soar:off / inhibit / watchdog)
  \item \texttt{landing\_gear} (on state change)
  \item \texttt{flight\_phase\_estimation}
\end{itemize}

%=============================================================================
\section{Run() --- Main Execution Loop}
%=============================================================================

\texttt{Run()} is the entry point called by the scheduler. Its complete
execution order each cycle:

\begin{enumerate}
  \item \textbf{Exit check:} if \texttt{should\_exit()}, unregister callback
        and clean up.
  \item \textbf{Parameter update:} if \texttt{parameter\_update} uORB updated,
        call \texttt{parameters\_update()}.
  \item \textbf{DDS soaring poll:} check \texttt{autosoaring\_control} uORB;
        apply inhibit/override/watchdog logic; optionally republish
        \texttt{autosoaring\_status} at 1\,Hz while the post-\texttt{soar:off}
        DDS inhibit is active (see Section~\ref{sec:dds}).
  \item \textbf{Global position:} update $(\phi, \lambda)$ from
        \texttt{vehicle\_global\_position}.
  \item \textbf{Altitude reference:} $h = -z_{\text{NED}} + h_{\text{ref}}$
        where $h_{\text{ref}}$ is the EKF reference altitude.
  \item \textbf{EKF reset handling} (manual modes only):
        \begin{itemize}
          \item Altitude reset: call \texttt{\_tecs.handle\_alt\_step()} to
                prevent a TECS altitude-error spike.
          \item XY reset: clear heading-hold flag.
        \end{itemize}
  \item \textbf{Setpoint ingestion:}
        \begin{itemize}
          \item Offboard (\texttt{flag\_control\_offboard\_enabled}):
                convert \texttt{trajectory\_setpoint} NED position/velocity
                into \texttt{position\_setpoint\_triplet} in lat/lon/alt.
                If velocity + acceleration are valid, compute signed
                loiter radius from lateral acceleration:
                \begin{equation}
                  R = -\frac{v_x^2 + v_y^2}{\|a_\perp\|}
                  \cdot \text{sign}(\hat{\mathbf{v}}_H \times \hat{\mathbf{a}}_\perp)
                \end{equation}
          \item Normal: subscribe \texttt{position\_setpoint\_triplet};
                set validity flags for prev/current/next.
                Reset altitude FOH counter \texttt{\_min\_current\_sp\_distance\_xy}.
        \end{itemize}
  \item \textbf{Sensor polls} (in order):
        \texttt{airspeed\_poll()}, \texttt{manual\_control\_setpoint\_poll()},
        \texttt{vehicle\_attitude\_poll()}, \texttt{vehicle\_command\_poll()},
        \texttt{vehicle\_control\_mode\_poll()}, \texttt{wind\_poll()}.
  \item \textbf{Air data:} update air density $\rho$ from
        \texttt{vehicle\_air\_data}; push density-dependent TECS parameters:
        \begin{equation}
          \dot{h}_{\text{climb,max}} = f(\rho), \quad
          \dot{h}_{\text{sink,min}} = f(\rho)
        \end{equation}
  \item \textbf{Land detection:} update \texttt{\_landed}.
  \item \textbf{Vehicle status:} update; clear VTOL backtransition state if
        not in transition.
  \item \textbf{Ground speed vector:} $(v_N, v_E)$ from local position.
  \item \textbf{Mode selection:} \texttt{set\_control\_mode\_current()}.
  \item \textbf{In-air states:} \texttt{update\_in\_air\_states()}.
  \item \textbf{TECS speed weight restore:} write \texttt{FW\_T\_SPDWEIGHT}
        back to TECS (may have been overridden in landing).
  \item \textbf{NPFG period restore:} write \texttt{NPFG\_PERIOD} back
        (may have been overridden in takeoff).
  \item \textbf{Default outputs:} flaps = 0, spoilers = 0, gear = keep,
        \texttt{fw\_control\_yaw\_wheel = false},
        \texttt{reset\_integral = false}.
  \item \textbf{Landing/takeoff state reset} (if not in respective mode).
  \item \textbf{Mode dispatch switch} (17 cases --- see Section~\ref{sec:dispatch}).
  \item \textbf{Post-dispatch:} throttle slew-rate gate; \texttt{status\_publish()};
        \texttt{landing\_status\_publish()}; gear publication; actuator
        setpoint publish; TECS debug publish.
\end{enumerate}

%=============================================================================
\section{Sensor Poll Functions}
%=============================================================================

\subsection{airspeed\_poll()}

\begin{enumerate}
  \item If \texttt{FW\_USE\_AIRSPD} enabled and \texttt{airspeed\_validated}
        has a new message:
        \begin{itemize}
          \item Check both \texttt{calibrated\_airspeed\_m\_s} and
                \texttt{true\_airspeed\_m\_s} are finite.
          \item Set \texttt{\_airspeed\_eas} (EAS) and compute
                EAS-to-TAS factor:
                \begin{equation}
                  k = \text{clamp}\!\left(
                    \frac{V_{\text{TAS,meas}}}{V_{\text{EAS,meas}}},\;0.9,\;2.0
                  \right)
                \end{equation}
          \item The clamp prevents out-of-range ISA atmosphere states.
          \item Record \texttt{\_time\_airspeed\_last\_valid}.
        \end{itemize}
  \item If no update for $> 1\,\text{s}$: mark airspeed invalid.
  \item On validity change: call \texttt{\_tecs.enable\_airspeed(valid)}
        to switch TECS between airspeed and no-airspeed control paths.
\end{enumerate}

\subsection{wind\_poll()}

\begin{enumerate}
  \item If \texttt{wind} uORB updated: store $(w_N, w_E)$; set valid if
        both components are finite.
  \item Timeout: if last update $> \texttt{WIND\_EST\_TIMEOUT}$, invalidate.
  \item If wind invalid: zero $(w_N, w_E)$ to prevent stale wind from
        contaminating NPFG (which uses wind for track-keeping).
\end{enumerate}

\subsection{manual\_control\_setpoint\_poll()}

\begin{enumerate}
  \item Copy latest stick input from \texttt{manual\_control\_setpoint}.
  \item Normal mapping:
        \begin{equation}
          \text{height\_rate\_stick} = \delta_{\text{pitch}}, \qquad
          \text{airspeed\_stick} = \delta_{\text{throttle}}
        \end{equation}
  \item Alternate mapping (\texttt{FW\_POS\_STK\_CONF} bit 0 set):
        \begin{equation}
          \text{height\_rate\_stick} = -\delta_{\text{throttle}}, \qquad
          \text{airspeed\_stick} = \delta_{\text{pitch}}
        \end{equation}
  \item Timeout: if last update $> 1\,\text{s}$, zero both to prevent
        stale input integrating into height or airspeed.
\end{enumerate}

\subsection{vehicle\_attitude\_poll()}

\begin{enumerate}
  \item Read quaternion $\mathbf{q}$ from \texttt{vehicle\_attitude};
        read body rates from \texttt{vehicle\_angular\_velocity}.
  \item Tailsitter correction: if VTOL tailsitter, rotate DCM by
        $90^\circ$ pitch offset before extracting Euler angles.
  \item Extract $(\phi, \theta, \psi)$ using \texttt{Eulerf(Dcmf(q))}.
  \item Compute body-frame acceleration:
        $\mathbf{a}_b = R^\top (a_N, a_E, a_D)^\top$.
        Store $\mathbf{a}_{b,x}$ for ground-speed undershoot detection.
  \item Compute body-frame velocity: store $\mathbf{v}_{b,x}$.
  \item Update TECS load factor:
        \begin{equation}
          n_z = \frac{1}{\cos\phi}, \qquad
          \texttt{\_tecs.set\_load\_factor}(n_z)
        \end{equation}
        This feeds the banked-stall speed floor (see Report 2, M9).
\end{enumerate}

\subsection{vehicle\_control\_mode\_poll()}

\begin{enumerate}
  \item If updated: check if arming state just changed
        (was \texttt{flag\_armed = false}, now true).
  \item On arming: call \texttt{reset\_takeoff\_state()} and
        \texttt{reset\_landing\_state()} to flush stale integrator states.
\end{enumerate}

\subsection{vehicle\_command\_poll()}

This function processes vehicle commands from uORB in a while-loop (multiple
commands may arrive between cycles).

\paragraph{DO\_GO\_AROUND}
Aborts landing if in \texttt{FW\_POSCTRL\_MODE\_AUTO} with
\texttt{SETPOINT\_TYPE\_LAND} and position control is enabled.
Calls \texttt{updateLandingAbortStatus(ABORTED\_BY\_OPERATOR)}, which
sets a flag that makes \texttt{control\_auto\_loiter()} fly level
until above a clearance altitude buffer.

\paragraph{DO\_CHANGE\_SPEED}
If \texttt{param1 = SPEED\_TYPE\_AIRSPEED} and \texttt{param2 > 0}:
\begin{itemize}
  \item Auto mode: overrides \texttt{pos\_sp\_triplet.current.cruising\_speed}
        (used by \texttt{adapt\_airspeed\_setpoint}).
  \item Manual mode: stores \texttt{\_commanded\_manual\_airspeed\_setpoint}
        (used by \texttt{get\_manual\_airspeed\_setpoint}).
\end{itemize}

\paragraph{DO\_SET\_MODE}
In the original system: no specific handling in
\texttt{FixedwingPositionControl} (handled by Commander).

%=============================================================================
\section{parameters\_update()}
%=============================================================================

Called on every \texttt{parameter\_update} event. Pushes all parameter
values into submodule instances.

\subsection{NPFG Parameters}

\begin{center}
\begin{tabular}{ll}
\toprule
PX4 Parameter & NPFG Method \\
\midrule
\texttt{NPFG\_PERIOD}            & \texttt{setPeriod()} \\
\texttt{NPFG\_DAMPING}           & \texttt{setDamping()} \\
\texttt{NPFG\_EN\_PERIOD\_LB/UB} & \texttt{enablePeriodLB/UB()} \\
\texttt{NPFG\_EN\_MIN\_GSP}      & \texttt{enableMinGroundSpeed()} \\
\texttt{NPFG\_EN\_TRACK\_KEEPING}& \texttt{enableTrackKeeping()} \\
\texttt{NPFG\_EN\_WIND\_REG}     & \texttt{enableWindExcessRegulation()} \\
\texttt{FW\_GND\_SPD\_MIN}       & \texttt{setMinGroundSpeed()} \\
\texttt{NPFG\_ROLL\_TC}          & \texttt{setRollTimeConst()} \\
\texttt{FW\_R\_LIM}              & \texttt{setRollLimit()} (converted to rad) \\
\texttt{NPFG\_SW\_DST\_MULT}     & \texttt{setSwitchDistanceMultiplier()} \\
\texttt{NPFG\_PERIOD\_SF}        & \texttt{setPeriodSafetyFactor()} \\
\bottomrule
\end{tabular}
\end{center}

\subsection{TECS Parameters}

\begin{center}
\begin{tabular}{ll}
\toprule
PX4 Parameter & TECS Method \\
\midrule
\texttt{FW\_T\_CLMB\_MAX} (via perf model, $\rho$) & \texttt{set\_max\_climb\_rate()} \\
\texttt{FW\_T\_SINK\_MAX}         & \texttt{set\_max\_sink\_rate()} \\
min sink rate (via perf model, $\rho$) & \texttt{set\_min\_sink\_rate()} \\
\texttt{FW\_T\_SPDWEIGHT}         & \texttt{set\_speed\_weight()} \\
trim/min/max EAS (via perf model) & \texttt{set\_equivalent\_airspeed\_*} \\
\texttt{FW\_T\_THR\_DAMPING}      & \texttt{set\_throttle\_damp()} \\
\texttt{FW\_T\_THR\_INTEG}        & \texttt{set\_integrator\_gain\_throttle()} \\
\texttt{FW\_T\_I\_GAIN\_PIT}      & \texttt{set\_integrator\_gain\_pitch()} \\
\texttt{FW\_THR\_SLEW\_MAX}       & \texttt{set\_throttle\_slewrate()} \\
\texttt{FW\_T\_VERT\_ACC}         & \texttt{set\_vertical\_accel\_limit()} \\
\texttt{FW\_T\_RLL2THR}           & \texttt{set\_roll\_throttle\_compensation()} \\
\texttt{FW\_T\_PTCH\_DAMP}        & \texttt{set\_pitch\_damping()} \\
\texttt{FW\_T\_H\_ERROR\_TC}      & \texttt{set\_altitude\_error\_time\_constant()} \\
\texttt{FW\_T\_FAST\_ALT\_ERR}    & \texttt{set\_fast\_descend\_altitude\_error()} \\
\texttt{FW\_T\_HRATE\_FF}         & \texttt{set\_altitude\_rate\_ff()} \\
\texttt{FW\_T\_TAS\_ERROR\_TC}    & \texttt{set\_airspeed\_error\_time\_constant()} \\
\texttt{STE\_RATE\_TC}            & \texttt{set\_ste\_rate\_time\_const()} \\
\texttt{SEB\_RATE\_FF}            & \texttt{set\_seb\_rate\_ff\_gain()} \\
\texttt{ASPD\_FMT\_STD\_DEV}      & \texttt{set\_airspeed\_measurement\_std\_dev()} \\
\texttt{ASPD\_RATE\_STD\_DEV}     & \texttt{set\_airspeed\_rate\_measurement\_std\_dev()} \\
\texttt{ASPD\_PROC\_NOISE}        & \texttt{set\_airspeed\_filter\_process\_std\_dev()} \\
\bottomrule
\end{tabular}
\end{center}

The performance model (\texttt{FixedwingPerformanceModel}) computes
density-dependent limits from the aerodynamic polars and \texttt{FW\_T\_CLMB\_MAX},
\texttt{FW\_T\_SINK\_MIN}. On every \texttt{parameters\_update()} call,
\texttt{\_performance\_model.runSanityChecks()} validates that min $<$ trim
$<$ max airspeed.

%=============================================================================
\section{set\_control\_mode\_current(): Mode Decision Tree}
\label{sec:setmode}
%=============================================================================

This function runs every cycle \emph{after} all polls, selecting one of
eight control modes:

\begin{center}
\begin{tabular}{lp{8cm}}
\toprule
Mode & Condition \\
\midrule
\texttt{FW\_POSCTRL\_MODE\_OTHER} & Rotary-wing vehicle type without transition, or no matching condition. \\
\texttt{FW\_POSCTRL\_MODE\_AUTO\_PATH} & Offboard enabled + position enabled + finite NED velocity setpoint. \\
\texttt{FW\_POSCTRL\_MODE\_AUTO} & Auto + position enabled + valid current setpoint (or IDLE type). \\
\texttt{FW\_POSCTRL\_MODE\_AUTO\_TAKEOFF} & As above, but setpoint type is TAKEOFF (fixed-wing only). \\
\texttt{FW\_POSCTRL\_MODE\_AUTO\_LANDING\_STRAIGHT} & As above, type LAND + prev setpoint valid. \\
\texttt{FW\_POSCTRL\_MODE\_AUTO\_LANDING\_CIRCULAR} & As above, type LAND + prev not valid. \\
\texttt{FW\_POSCTRL\_MODE\_AUTO\_ALTITUDE} & Auto + climb-rate enabled + armed + position \emph{not} valid; GPS failsafe phase 1 (fixed-bank loiter for \texttt{NAV\_GPSF\_LT} seconds). \\
\texttt{FW\_POSCTRL\_MODE\_AUTO\_CLIMBRATE} & Same conditions, after \texttt{NAV\_GPSF\_LT} timeout; GPS failsafe phase 2 (emergency descent). \\
\texttt{FW\_POSCTRL\_MODE\_MANUAL\_POSITION} & Manual + position enabled. \\
\texttt{FW\_POSCTRL\_MODE\_MANUAL\_ALTITUDE} & Manual + altitude enabled. \\
\texttt{FW\_POSCTRL\_MODE\_TRANSITION\_TO\_HOVER\_*} & VTOL back-transition. \\
\bottomrule
\end{tabular}
\end{center}

\textbf{VTOL takeoff bypass:} If VTOL and in transition while setpoint
type is TAKEOFF, the type is rewritten to \texttt{POSITION} and mode is
set to \texttt{FW\_POSCTRL\_MODE\_AUTO} to use the standard waypoint
follower instead of the launch sequence.

\textbf{Arming guard:} Failsafe altitude/climbrate modes only activate
when the vehicle is armed (\texttt{flag\_armed = true}) to prevent triggering
during pre-arm checks.

%=============================================================================
\section{Run() Mode Dispatch Switch}
\label{sec:dispatch}
%=============================================================================

After \texttt{set\_control\_mode\_current()}, a switch dispatches to the
appropriate control function:

\begin{lstlisting}[language=C++]
switch (_control_mode_current) {
  case FW_POSCTRL_MODE_AUTO:
      control_auto(...);                     break;
  case FW_POSCTRL_MODE_AUTO_ALTITUDE:
      control_auto_fixed_bank_alt_hold(...); break;
  case FW_POSCTRL_MODE_AUTO_CLIMBRATE:
      control_auto_descend(...);             break;
  case FW_POSCTRL_MODE_AUTO_LANDING_STRAIGHT:
      control_auto_landing_straight(...);    break;
  case FW_POSCTRL_MODE_AUTO_LANDING_CIRCULAR:
      control_auto_landing_circular(...);    break;
  case FW_POSCTRL_MODE_AUTO_PATH:
      control_auto_path(...);                break;
  case FW_POSCTRL_MODE_AUTO_TAKEOFF:
      control_auto_takeoff(...);             break;
  case FW_POSCTRL_MODE_MANUAL_POSITION:
      control_manual_position(...);          break;
  case FW_POSCTRL_MODE_MANUAL_ALTITUDE:
      control_manual_altitude(...);          break;
  case FW_POSCTRL_MODE_TRANSITION_TO_HOVER_LINE_FOLLOW:
      control_backtransition_line_follow(...); break;
  case FW_POSCTRL_MODE_TRANSITION_TO_HOVER_HEADING_HOLD:
      control_backtransition_heading_hold(); break;
}
\end{lstlisting}

%=============================================================================
\section{NPFG: Nonlinear Path Following Guidance}
%=============================================================================

\subsection{Core Guidance Law}

NPFG computes a lateral acceleration command based on the error between
the current airspeed vector and the desired track direction.

\paragraph{Track-error angle $\eta$:}
For a line segment $A \to B$ with unit tangent $\hat{\mathbf{t}}$:
\begin{equation}
  e_y = \hat{\mathbf{n}} \cdot (\mathbf{p} - \mathbf{A})
  \quad\text{(signed cross-track error)}
\end{equation}

\begin{equation}
  \eta = \text{atan2}(\mathbf{v}_a, \hat{\mathbf{t}}) + \arcsin\!\left(
    \text{clamp}\!\left(\frac{-e_y}{L},\;-1,\;1\right)
  \right)
\end{equation}

where $L$ is the look-ahead distance.

\paragraph{Look-ahead length:}
\begin{equation}
  L = \frac{T_{\text{NPFG}}\,V_a}{2\pi / \zeta} = f(\text{NPFG period, damping})
\end{equation}

Bounded below by a minimum ground speed term and above by a maximum.

\paragraph{Lateral acceleration command:}
\begin{equation}
  a_{\text{lat}} = \frac{2V_a^2}{L}\,\sin\eta
\end{equation}

\paragraph{For loiter circle:} The centripetal acceleration is added:
\begin{equation}
  a_{\text{lat}} \mathrel{+}= \text{sign}(d) \cdot \frac{V_a^2}{R}
\end{equation}

where $d = \pm 1$ for CW/CCW.

\paragraph{Roll setpoint:}
\begin{equation}
  \phi^* = \arctan\!\left(\frac{a_{\text{lat}}}{g}\right)
\end{equation}

Clamped to $[-\phi_{\max}, \phi_{\max}]$ from \texttt{FW\_R\_LIM}.

\subsection{getCorrectedNpfgRollSetpoint()}

NPFG reports a ``can-run factor'' $f_{\text{run}} \in [0, 1]$ that
measures confidence based on:
\begin{itemize}
  \item Velocity estimate validity (\texttt{xy\_valid} in local position).
  \item Wind estimate validity (\texttt{\_wind\_valid}).
\end{itemize}

The corrected roll:
\begin{equation}
  \phi^*_{\text{corrected}} = f_{\text{run}} \cdot \phi^*_{\text{NPFG}}
\end{equation}

If $f_{\text{run}} < 0.9$ for more than \texttt{ROLL\_WARNING\_TIMEOUT}
(2\,s), a user warning event is emitted. This prevents the aircraft from
making large rolls in GPS-degraded conditions.

\textbf{Roll slew rate:} After NPFG, the roll setpoint is filtered
through a slew-rate limiter (\texttt{FW\_PN\_R\_SLEW\_MAX}, default
90\,deg/s) to prevent step demands.

%=============================================================================
\section{adapt\_airspeed\_setpoint()}
%=============================================================================

Called before every TECS invocation. Six processing steps:

\begin{enumerate}
  \item \textbf{Wind-based minimum airspeed scaling:}
        If airspeed valid, wind valid, and \texttt{FW\_WIND\_ARSP\_SC} $> 0$:
        \begin{equation}
          V_{\text{EAS,min}} \leftarrow \min(V_{\text{EAS,min}} +
            k_w\,\|\mathbf{w}\|,\; V_{\text{EAS,max}})
        \end{equation}
        This increases the minimum airspeed in strong winds to maintain
        climb performance and avoid stall.

  \item \textbf{NaN/zero setpoint fallback:}
        If the mission waypoint cruising speed is not finite or $\leq 0$:
        $V^*_{\text{EAS}} \leftarrow V_{\text{trim}}$.

  \item \textbf{Ground-speed undershoot correction:}
        If wind is \emph{not} valid (cannot estimate it):
        \begin{equation}
          \text{if } v_{g,x} < V_{g,\min}: \quad
          V^*_{\text{EAS}} \mathrel{+}= V_{g,\min} - v_{g,x}
        \end{equation}
        where $v_{g,x}$ is the body-frame forward ground speed. This
        prevents the aircraft being blown backwards.

  \item \textbf{Bounds:}
        $V^*_{\text{EAS}} = \text{clamp}(V^*_{\text{EAS}}, V_{\text{EAS,min}}, V_{\text{EAS,max}})$.

  \item \textbf{Slew-rate controller initialisation:}
        On first call or after TECS reset:
        $V^*_{\text{slew}} = \max(V_{\text{EAS,min}}, V_{\text{EAS,meas}})$.
        If the current state is below $V_{\text{EAS,min}}$ (e.g.\ minimum
        changed), force the slew state up to the new minimum.

  \item \textbf{Slew-rate update:}
        $V^*_{\text{out}} = \texttt{SlewRate.update}(V^*_{\text{EAS}}, \Delta t)$
        with rate limit \texttt{ASPD\_SP\_SLEW\_RATE} (default $1\,\text{m/s}^2$).
\end{enumerate}

%=============================================================================
\section{get\_manual\_airspeed\_setpoint()}
%=============================================================================

Converts the throttle stick position $\delta_T \in [-1, 1]$ to an EAS
setpoint for manual modes:

\begin{itemize}
  \item \textbf{Manual airspeed enabled} (\texttt{FW\_POS\_STK\_CONF} bit 1):
        piecewise linear interpolation:
        \begin{equation}
          V^* = \texttt{interpolateNXY}\!\left(\delta_T,\;
            \{-1, 0, 1\},\;\{V_{\min}, V_{\text{trim}}, V_{\max}\}\right)
        \end{equation}
        Neutral stick ($\delta_T = 0$) gives trim airspeed.

  \item \textbf{Airspeed override} (\texttt{DO\_CHANGE\_SPEED} received):
        use the commanded value clamped to $[V_{\min}, V_{\max}]$.

  \item \textbf{Default}: return trim airspeed.
\end{itemize}

%=============================================================================
\section{getManualHeightRateSetpoint()}
%=============================================================================

Converts pitch-stick position to a height rate setpoint for manual altitude
control:

\begin{equation}
  \dot{h}^* = \begin{cases}
    \delta_{\text{pitch}} \cdot \dot{h}_{\text{climb,target}} & \delta_{\text{pitch}} > 0 \\
    -\delta_{\text{pitch}} \cdot \dot{h}_{\text{sink,target}} & \delta_{\text{pitch}} < 0 \\
    \text{NaN} & |\delta_{\text{pitch}}| < \varepsilon_{\text{dead}}
  \end{cases}
\end{equation}

NaN output means ``no height-rate control'' (TECS tracks altitude reference).

%=============================================================================
\section{tecs\_update\_pitch\_throttle()}
%=============================================================================

Single call site for TECS. Called by every position control sub-mode.

\textbf{Step 1: Inject limits into TECS.}
\begin{lstlisting}[language=C++]
_tecs.set_min_pitch(pitch_min_rad);
_tecs.set_max_pitch(pitch_max_rad);
_tecs.set_throttle_min(thr_min);
_tecs.set_throttle_max(thr_max);
_tecs.set_target_climbrate(target_climbrate);
_tecs.set_target_sinkrate(target_sinkrate);
\end{lstlisting}

\textbf{Step 2: Low-height mode.}
If \texttt{is\_low\_height} (aircraft below \texttt{FW\_LH\_APC\_FACTOR}
$\times$ landing approach clearance altitude), TECS uses tighter altitude
error gains and smaller pitch limits to improve precision near the ground.

\textbf{Step 3: Underspeed handling.}
If \texttt{disable\_underspeed\_handling = false} (default), the
\texttt{detect\_underspeed\_enabled} flag is true in TECS.

\textbf{Step 4: Height-rate override.}
If the caller provides a finite \texttt{height\_rate\_setpoint}, it
is passed to the TECS reference model as a direct height-rate command
(bypassing the altitude proportional controller).

\textbf{Step 5: EAS → TAS setpoint.}
TECS receives the raw EAS setpoint; internally calls
\texttt{calcTrueAirspeedSetpoint()} (see Report 2).

\textbf{Step 6: Call TECS::update()}.

\textbf{Step 7: Output.}
\begin{itemize}
  \item \texttt{get\_tecs\_pitch()} returns the TECS pitch setpoint $\theta^*$.
  \item \texttt{get\_tecs\_thrust()} returns the TECS throttle $\delta_t^*$
        clamped to $[\delta_{t,\min}, \delta_{t,\max}]$.
\end{itemize}

%=============================================================================
\section{control\_auto(): Setpoint Type Dispatch}
%=============================================================================

\texttt{control\_auto()} runs when in \texttt{FW\_POSCTRL\_MODE\_AUTO}.

\subsection{handle\_setpoint\_type()}

Before the inner switch, the setpoint type is re-evaluated:
\begin{itemize}
  \item If velocity control only (no position): return
        \texttt{SETPOINT\_TYPE\_VELOCITY}.
  \item If type is \texttt{POSITION} and not approaching VTOL
        backtransition:
        compute 3D distance to waypoint. If aircraft is within the
        acceptance radius AND has the correct vertical separation:
        \begin{equation}
          d_{xy} < r_{\text{acc}}, \quad d_z > r_{z,\text{acc}}
          \;\Rightarrow\; \text{switch to LOITER}
        \end{equation}
        $r_{\text{acc}}$ comes from \texttt{\_npfg.switchDistance(500)}
        which scales with airspeed.
\end{itemize}

\subsection{Inner Switch in control\_auto()}

\begin{lstlisting}[language=C++]
switch (position_sp_type) {
  case SETPOINT_TYPE_IDLE:    thrust = 0; break;
  case SETPOINT_TYPE_POSITION:
      control_auto_position(...); break;
  case SETPOINT_TYPE_VELOCITY:
      control_auto_velocity(...); break;
  case SETPOINT_TYPE_LOITER:
      control_auto_loiter(...);   break;
#ifdef CONFIG_FIGURE_OF_EIGHT
  case SETPOINT_TYPE_FIGURE_EIGHT:
      controlAutoFigureEight(...); break;
#endif
}
\end{lstlisting}

%=============================================================================
\section{control\_auto\_position()}
%=============================================================================

\begin{enumerate}
  \item \textbf{Altitude First-Order Hold (FOH):}
        Linearly interpolates the altitude setpoint between previous
        and current waypoints as the aircraft approaches:
        \begin{equation}
          h^* = h_{\text{prev}} +
            \frac{h_{\text{curr}} - h_{\text{prev}}}{d_{A \to B} - r_{\text{acc}}}
            \cdot (d_{A \to B} - d_{\text{min}})
        \end{equation}
        where $d_{\text{min}}$ is the smallest distance to the current
        waypoint ever achieved. This ensures the aircraft reaches the
        waypoint altitude precisely at the acceptance radius rather than
        demanding a step change.

  \item \textbf{NPFG guidance:}
        Calls \texttt{navigateWaypoints(prev, curr, ...)} if a valid
        previous waypoint exists; else \texttt{navigateWaypoint(curr,
        ...)}. Outputs roll setpoint $\phi^*$.

  \item \textbf{Airspeed setpoint:}
        \texttt{adapt\_airspeed\_setpoint(pos\_sp\_curr.cruising\_speed, \ldots)}.
        After NPFG: $V^* \leftarrow$ \texttt{getAirspeedRef()} / $k$
        (NPFG may increase the airspeed demand to maintain track in
        crosswinds).

  \item \textbf{TECS call:} \texttt{tecs\_update\_pitch\_throttle}$(h^*, V^*, \ldots)$.

  \item \textbf{Attitude setpoint:}
        $\mathbf{q}^* = \text{Euler}(\phi^*_{\text{NPFG}}, \theta^*_{\text{TECS}}, \psi)$.
\end{enumerate}

%=============================================================================
\section{control\_auto\_loiter()}
%=============================================================================

\begin{enumerate}
  \item \textbf{Radius:} from \texttt{pos\_sp\_curr.loiter\_radius}; if zero,
        use \texttt{NAV\_LOITER\_RAD} (default $50\,\text{m}$).

  \item \textbf{Early landing configuration:}
        If the next setpoint is type LAND and the aircraft is within
        $R + r_{\text{switch}}$ of the loiter circle centre,
        and \texttt{FW\_LND\_EARLYCFG} enabled:
        \begin{itemize}
          \item Set airspeed to landing airspeed \texttt{FW\_LND\_AIRSPD}
                (or stall margin if not set).
          \item Deploy flaps: \texttt{\_flaps\_setpoint = FW\_FLAPS\_LND\_SCL}.
          \item Deploy spoilers: \texttt{\_spoilers\_setpoint = FW\_SPOILERS\_LND}.
          \item Lower gear: \texttt{GEAR\_DOWN}.
          \item Enable low-height TECS mode.
        \end{itemize}

  \item \textbf{Landing abort override:}
        If \texttt{\_landing\_abort\_status} is set:
        \begin{itemize}
          \item If altitude delta to loiter $< \texttt{kClearanceAltitudeBuffer}$
                (10\,m): clear abort status (done).
          \item Else: set roll to 0 (fly straight), retract flaps if
                airspeed is safe.
        \end{itemize}

  \item \textbf{NPFG:} \texttt{navigateLoiter(centre, pos, R, direction, ...)}
        uses the loiter-circle guidance.

  \item \textbf{TECS call:} with loiter altitude setpoint.

  \item \textbf{Attitude setpoint:}
        $\mathbf{q}^* = \text{Euler}(\phi^*_{\text{NPFG}}, \theta^*_{\text{TECS}}, \psi)$.
\end{enumerate}

%=============================================================================
\section{control\_auto\_velocity()}
%=============================================================================

Used for offboard velocity setpoints (no position, only velocity).

\begin{enumerate}
  \item Target velocity from \texttt{pos\_sp\_curr.vx/vy};
        heading from \texttt{atan2(vy, vx)}.
  \item NPFG navigates toward a virtual waypoint in the target direction.
  \item TECS controls altitude to \texttt{pos\_sp\_curr.alt}.
\end{enumerate}

%=============================================================================
\section{control\_auto\_path()}
%=============================================================================

Used for offboard path following with position + velocity + acceleration.

\begin{enumerate}
  \item If acceleration is available: compute lateral curvature from
        the cross-product of normalised velocity and acceleration;
        use signed loiter radius from centripetal acceleration.
  \item NPFG path follows the straight segment defined by
        velocity direction.
  \item TECS controls altitude.
\end{enumerate}

%=============================================================================
\section{control\_auto\_fixed\_bank\_alt\_hold() --- GPS Failsafe Phase 1}
%=============================================================================

Activated when the GPS position estimate is lost
(\texttt{FW\_POSCTRL\_MODE\_AUTO\_ALTITUDE}) for up to \texttt{NAV\_GPSF\_LT}
seconds (default 30\,s).

\begin{itemize}
  \item Roll set to \texttt{NAV\_GPSF\_R} (fixed open-loop bank angle,
        default $15^\circ$): aircraft circles indefinitely.
  \item Yaw: 0 (heading holds whatever it was).
  \item TECS controls altitude to current measured altitude, trim airspeed.
  \item Special: if NED $z$ or $v_z$ is invalid (no barometer):
        thrust forced to \texttt{FW\_THR\_MIN} to avoid climb-away.
\end{itemize}

%=============================================================================
\section{control\_auto\_descend() --- GPS Failsafe Phase 2}
%=============================================================================

Activated after \texttt{NAV\_GPSF\_LT} timeout.

\begin{itemize}
  \item Height rate setpoint hard-coded to $-0.5\,\text{m/s}$.
  \item Same fixed-bank open-loop roll as Phase 1.
  \item TECS passed this height rate directly (bypasses altitude controller).
  \item If $v_z$ invalid: thrust forced to minimum.
\end{itemize}

%=============================================================================
\section{control\_auto\_takeoff()}
%=============================================================================

Multi-phase launch sequence (catapult, runway, VTOL front-transition):

\begin{enumerate}
  \item \textbf{Pre-launch:} no airspeed, no attitude control; wait for
        acceleration signature or RC trigger.
  \item \textbf{Climbout:} TECS active, pitch clamped above
        \texttt{FW\_TKO\_PITCH\_MIN} to ensure ground clearance.
        NPFG follows the runway heading (or waypoint).
  \item \textbf{Transition altitude reached:} gear up; switch to normal
        auto-position mode.
\end{enumerate}

%=============================================================================
\section{control\_auto\_landing\_straight() and \_circular()}
%=============================================================================

\subsection{control\_auto\_landing\_straight() --- Mission Landing}

\begin{enumerate}
  \item Uses previous WP as the glideslope start; current WP as
        touchdown.
  \item Glideslope: altitude setpoint linearly interpolated from approach
        altitude to runway elevation.
  \item Flare: when within \texttt{FW\_LND\_FLALT} of touchdown:
        \begin{itemize}
          \item Throttle cut (\texttt{FW\_THR\_MIN}).
          \item Pitch held at \texttt{FW\_LND\_FL\_PMAX}.
          \item Airspeed setpoint switches to \texttt{FW\_LND\_AIRSPD}.
        \end{itemize}
\end{enumerate}

\subsection{control\_auto\_landing\_circular() --- Hold Pattern Landing}

Used when no previous waypoint is available (e.g.\ Return-to-Land mode).
The aircraft spirals down in a loiter until below the decision altitude,
then transitions to straight final.

%=============================================================================
\section{control\_manual\_altitude()}
%=============================================================================

\begin{enumerate}
  \item \texttt{updateManualTakeoffStatus()}: detect first climb above
        takeoff threshold; set \texttt{\_completed\_manual\_takeoff}.
  \item Airspeed setpoint: \texttt{adapt\_airspeed\_setpoint(get\_manual\_airspeed\_setpoint(), \ldots)}.
  \item Height rate: \texttt{getManualHeightRateSetpoint()} from pitch stick.
  \item Pitch minimum: \texttt{FW\_P\_LIM\_MIN} normally; or
        \texttt{MIN\_PITCH\_DURING\_MANUAL\_TAKEOFF} ($8^\circ$) if
        not yet completed takeoff.
  \item Throttle maximum: if on ground and stick below
        \texttt{THROTTLE\_THRESH} (0.05): $\delta_{t,\max} = 0$ (kill switch).
  \item Roll: direct from stick scaled by \texttt{FW\_R\_LIM}.
  \item TECS + attitude setpoint.
\end{enumerate}

%=============================================================================
\section{control\_manual\_position()}
%=============================================================================

Manual position mode adds heading hold on top of manual altitude:

\begin{enumerate}
  \item Same airspeed and height-rate processing as altitude mode.
  \item \textbf{Heading hold engagement condition:}
        Stick roll and yaw both below \texttt{HDG\_HOLD\_MAN\_INPUT\_THRESH}
        ($0.01$) AND yaw rate below \texttt{HDG\_HOLD\_YAWRATE\_THRESH}
        ($0.15\,\text{rad/s}$).
  \item When engaged:
        \begin{itemize}
          \item Record hold yaw $\psi_{\text{hold}}$ and hold position
                $\mathbf{p}_{\text{hold}}$.
          \item NPFG \texttt{navigateLine()} from hold position along
                heading $\psi_{\text{hold}}$.
          \item Roll from NPFG; yaw from $\psi_{\text{hold}}$.
        \end{itemize}
  \item EKF XY reset during heading hold: update hold position for 2\,s
        to let EKF converge.
  \item When stick moves: disable heading hold, use direct roll stick.
\end{enumerate}

%=============================================================================
\section{control\_backtransition\_heading\_hold() and \_line\_follow()}
%=============================================================================

Used for VTOL back-transition (forward flight $\to$ hover).

\subsection{heading\_hold}
\begin{itemize}
  \item Lock the heading at transition start.
  \item Simulate wind = 0; ground speed = airspeed aligned with heading.
  \item NPFG navigates a virtual waypoint 200\,m ahead.
  \item Thrust forced to \texttt{FW\_THR\_MIN} (overridden by transition
        logic in VTOL module).
\end{itemize}

\subsection{line\_follow}
\begin{itemize}
  \item NPFG follows the line from back-transition start position to
        land waypoint.
  \item Position at back-transition start captured on first call.
\end{itemize}

%=============================================================================
\section{Throttle Output (Original)}
%=============================================================================

After the mode dispatch, throttle is applied:

\begin{lstlisting}[language=C++]
if (position_sp_type == SETPOINT_TYPE_IDLE) {
    _att_sp.thrust_body[0] = 0.0f;
} else if (_landed && ...) {
    _att_sp.thrust_body[0] = min(_param_fw_thr_idle.get(), 1.f);
} else {
    _att_sp.thrust_body[0] = min(get_tecs_thrust(), _param_fw_thr_max.get());
}
\end{lstlisting}

No ramp logic. TECS output goes directly to the airframe actuators.

% ============================================================
%  PART II  — SOARING-MODE MODIFICATIONS
% ============================================================
\part{Soaring-Mode Modifications and Their Physical Impact}

%=============================================================================
\section{Overview of Modified Functions}
%=============================================================================

The soaring implementation modifies the following functions:

\begin{center}
\begin{tabular}{ll}
\toprule
Function & What changes \\
\midrule
\texttt{Run()} & DDS poll, inhibit/watchdog, status republish, DDS→cmd arbiter \\
\texttt{vehicle\_command\_poll()} & CUSTOM\_0, DO\_SET\_MODE, \texttt{publish\_autosoaring\_status()} on soar:off \\
\texttt{publish\_autosoaring\_status()} & uORB \texttt{autosoaring\_status} (bridged to ROS\,2) \\
\texttt{parameters\_update()} & Glide airspeed eager pre-push \\
\texttt{control\_auto()} & Full soaring dispatch block \\
Throttle output & Symmetric entry/exit ramps \\
\bottomrule
\end{tabular}
\end{center}

%=============================================================================
\section{Change 1: Run() --- DDS Poll, Inhibit, and Staleness Watchdog}
\label{sec:dds}
%=============================================================================

\subsection{DDS Poll}

\begin{lstlisting}[language=C++]
autosoaring_control_s soaring_msg;
if (_autosoaring_control_sub.update(&soaring_msg)) {
    const hrt_abstime now = hrt_absolute_time();
    const bool dds_inhibited = (_soaring_dds_inhibit_until_us > 0
        && now < _soaring_dds_inhibit_until_us);
    const bool dds_wants_soaring =
        (soaring_msg.soaring_mode != autosoaring_control_s::SOARING_OFF);

    if (dds_inhibited && dds_wants_soaring) {
        // ignore: CLI soar:off inhibit active
    } else {
        if (!dds_wants_soaring) {
            _soaring_dds_inhibit_until_us = 0; // DDS OFF cancels inhibit
        }
        if (!_soaring_local_override) {
            const bool prev_thermal = (
                _autosoaring_control.soaring_mode == autosoaring_control_s::SOARING_THERMAL_LOITER ||
                _autosoaring_control.soaring_mode == autosoaring_control_s::SOARING_THERMAL_BANK);
            const bool new_thermal = (
                soaring_msg.soaring_mode == autosoaring_control_s::SOARING_THERMAL_LOITER ||
                soaring_msg.soaring_mode == autosoaring_control_s::SOARING_THERMAL_BANK);
            _autosoaring_control = soaring_msg;
            _autosoaring_last_recv_us = now;
            if (!prev_thermal && new_thermal) {
                _soaring_mode_cmd_last_us = 0;
                _soaring_last_lat = NAN;
                _soaring_last_lon = NAN;
            }
        }
    }
}
\end{lstlisting}

\paragraph{DDS inhibit:} When \texttt{soar:off} is issued from CLI,
the DDS path is blocked for 5\,s (\texttt{\_soaring\_dds\_inhibit\_until\_us
= now + 5s}). This prevents an autonomous companion from immediately
re-enabling soaring after the operator has manually stopped it.
If the companion publishes \texttt{soaring\_mode = OFF} while the inhibit is
active, the remaining cooldown is cleared immediately (\texttt{\_soaring\_dds\_inhibit\_until\_us = 0}).

\paragraph{Authority arbitration:} The \texttt{\_soaring\_local\_override}
flag gives CLI exclusive authority. When set, DDS messages are received
but do \emph{not} overwrite \texttt{\_autosoaring\_control}. CLI and DDS
cooperate: CLI sets the mode; DDS can update ancillary fields (loiter
position, bank angle) while CLI holds authority.

\subsection{Staleness Watchdog}

\begin{lstlisting}[language=C++]
if (!_soaring_local_override &&
    _autosoaring_last_recv_us > 0 &&
    (hrt_absolute_time() - _autosoaring_last_recv_us) > 2_s) {
    if (_autosoaring_control.soaring_mode != autosoaring_control_s::SOARING_OFF) {
        _autosoaring_control.soaring_mode = autosoaring_control_s::SOARING_OFF;
        publish_autosoaring_status(autosoaring_status_s::SOURCE_STALENESS_WATCHDOG,
            autosoaring_control_s::SOARING_OFF, 0);
    }
    _tecs.set_gliding_mode_enabled(false);
    _autosoaring_last_recv_us = 0;
}
\end{lstlisting}

\textbf{Physical motivation:} The companion computer may crash or lose
its DDS connection while in mid-glide. Without a watchdog, the aircraft
would remain engine-off indefinitely at risk of altitude loss. The 2\,s
timeout restores powered flight automatically.

\textbf{Not triggered by CLI:} \texttt{\_autosoaring\_last\_recv\_us = 0}
when CLI is the authority (set in \texttt{vehicle\_command\_poll()} for
\texttt{CUSTOM\_0} commands). The watchdog only fires if
\texttt{\_autosoaring\_last\_recv\_us > 0}, which is only set when DDS
delivered a message.

\subsection{AutosoaringStatus publication, 1\,Hz inhibit republish, and 2\,Hz phase heartbeat}

On CLI \texttt{soar:off} (\texttt{CUSTOM\_0} with mode~0), after TECS gliding is
disabled, \texttt{publish\_autosoaring\_status()} is called with
\texttt{SOURCE\_CLI\_SOARING\_OFF}, the effective \texttt{soaring\_mode}
(\texttt{OFF}), \texttt{dds\_inhibit\_until\_us} equal to
\texttt{hrt\_absolute\_time() + 5\,s}, and
\texttt{fmu\_soaring\_phase = SOARING\_PHASE\_CRUISE}. The timestamp
\texttt{\_autosoaring\_status\_repub\_us} is initialised so that, for the
duration of the inhibit window, the same status is republished at 1\,Hz
(companion may miss a single DDS sample):

\begin{lstlisting}[language=C++]
if (_soaring_dds_inhibit_until_us > 0 && status_now < _soaring_dds_inhibit_until_us) {
    if (status_now - _autosoaring_status_repub_us >= 1_s) {
        _autosoaring_status_repub_us = status_now;
        publish_autosoaring_status(SOURCE_CLI_SOARING_OFF,
            _autosoaring_control.soaring_mode,
            _soaring_dds_inhibit_until_us,
            SOARING_PHASE_CRUISE);
    }
} else {
    _autosoaring_status_repub_us = 0;
}
\end{lstlisting}

In addition, a 2\,Hz periodic heartbeat is published at the end of every
\texttt{control\_auto()} call (inside the position-control rate loop at
\(\approx 50\)\,Hz), derived from the current \texttt{effective\_glide}/
\texttt{effective\_thermal} booleans:

\begin{lstlisting}[language=C++]
if (phase_now - _autosoaring_status_periodic_us >= 500_ms) {
    _autosoaring_status_periodic_us = phase_now;
    uint8_t fmu_phase;
    if (effective_thermal)
        fmu_phase = PX4_ISFINITE(_soaring_last_lat)
                    ? SOARING_PHASE_THERMAL : SOARING_PHASE_SEARCHING;
    else if (effective_glide)
        fmu_phase = SOARING_PHASE_GLIDING;
    else
        fmu_phase = SOARING_PHASE_CRUISE;
    publish_autosoaring_status(SOURCE_PERIODIC,
        _autosoaring_control.soaring_mode,
        _soaring_dds_inhibit_until_us, fmu_phase);
}
\end{lstlisting}

This means \texttt{AutosoaringStatus} is \textbf{not} limited to the 5\,s
inhibit window. It is published continuously: every 500\,ms during any active
soaring, at 1\,Hz during the DDS inhibit cooldown, and as an immediate
one-shot on event (watchdog or \texttt{soar:off}).
Because uORB automatically logs every published topic, the phase sequence is
available in the ULog for post-flight analysis.

The uORB topic is listed under \texttt{publications} in
\texttt{dds\_topics.yaml} as \texttt{/fmu/out/autosoaring\_status}.

%=============================================================================
\section{Change 2: vehicle\_command\_poll() --- CUSTOM\_0 and DO\_SET\_MODE}
%=============================================================================

\subsection{VEHICLE\_CMD\_CUSTOM\_0 (Soaring CLI)}

\begin{lstlisting}[language=C++]
const uint8_t new_mode = constrain((int)cmd.param1,
    0, (int)SOARING_THERMAL_BANK);
_autosoaring_control.soaring_mode    = new_mode;
_autosoaring_control.loiter_radius_m =
    (cmd.param3 > FLT_EPSILON) ? cmd.param3 : NAN;
if (PX4_ISFINITE(cmd.param2) && cmd.param2 > FLT_EPSILON)
    _autosoaring_control.bank_angle_cmd = cmd.param2;
\end{lstlisting}

On enabling (\texttt{new\_mode != SOARING\_OFF}):
\begin{itemize}
  \item \texttt{\_soaring\_local\_override = true}
  \item Clear inhibit (\texttt{\_soaring\_dds\_inhibit\_until\_us = 0})
  \item Clear latch (\texttt{\_soaring\_forbidden\_latched = false})
  \item Disable watchdog (\texttt{\_autosoaring\_last\_recv\_us = 0})
  \item Set ``paired mode-switch'' flag (\texttt{\_soaring\_expect\_set\_mode = true})
  \item Reset DO\_REPOSITION debounce (\texttt{\_soaring\_mode\_cmd\_last\_us = 0},
        \texttt{\_soaring\_last\_lat = NAN})
\end{itemize}

On disabling (\texttt{new\_mode == SOARING\_OFF}):
\begin{itemize}
  \item Clear override and latch
  \item Block DDS for 5\,s (\texttt{\_soaring\_dds\_inhibit\_until\_us = now + 5s})
  \item Immediately call \texttt{\_tecs.set\_gliding\_mode\_enabled(false)}
  \item Publish \texttt{autosoaring\_status} (\texttt{SOURCE\_CLI\_SOARING\_OFF},
        inhibit deadline) and seed the 1\,Hz republish timer
\end{itemize}

\subsection{VEHICLE\_CMD\_DO\_SET\_MODE}

When CLI has override, the module must distinguish between:
\begin{enumerate}
  \item A DO\_SET\_MODE \emph{paired} with CUSTOM\_0 (i.e.\ \texttt{soar:glide}
        sends both in sequence): the mode switch IS part of soaring
        activation. Consume \texttt{\_soaring\_expect\_set\_mode} and
        keep soaring on.
  \item A standalone DO\_SET\_MODE (operator typed
        \texttt{commander mode auto:mission}): interpret as ``exit soaring,
        restore powered flight''.
\end{enumerate}

\begin{lstlisting}[language=C++]
if (_soaring_local_override) {
    if (_soaring_expect_set_mode) {
        _soaring_expect_set_mode = false; // consume pairing
    } else {
        // Check target mode: ignore QGC heartbeats (AUTO_MISSION/LOITER)
        const uint8_t main = (uint8_t)cmd.param2;
        const uint8_t sub  = (uint8_t)cmd.param3;
        const bool compat  =
            (main == PX4_CUSTOM_MAIN_MODE_AUTO &&
             (sub == PX4_CUSTOM_SUB_MODE_AUTO_MISSION ||
              sub == PX4_CUSTOM_SUB_MODE_AUTO_LOITER));
        if (!compat) {
            // Incompatible mode (MANUAL, STABILIZED, etc.) -> exit soaring
            _autosoaring_control.soaring_mode = SOARING_OFF;
            _soaring_local_override    = false;
            _tecs.set_gliding_mode_enabled(false);
        }
        // else: AUTO_MISSION / AUTO_LOITER heartbeat -> keep soaring
    }
}
\end{lstlisting}

\paragraph{QGC heartbeat fix.}
QGC periodically broadcasts \texttt{DO\_SET\_MODE(AUTO\_MISSION)} to maintain
the displayed mode. Without the \texttt{compat} guard, every such heartbeat
would arrive as a ``standalone'' mode-switch and immediately disable soaring —
a race condition that was observed in SITL testing.

The guard keeps soaring active when the target mode is \texttt{AUTO\_MISSION}
or \texttt{AUTO\_LOITER} (both are soaring-compatible). Any other mode
(MANUAL, ACRO, STABILIZED, ALTITUDE, POSCTL) is considered an intentional
operator override and correctly exits soaring.

\paragraph{Interaction with altitude floor.}
The QGC fix keeps \texttt{\_soaring\_local\_override = true} permanently during
CLI glide, which had the unintended side effect of also keeping the
\texttt{FW\_ALT\_MIN} check bypassed (the old code guarded it with
\texttt{!\_soaring\_local\_override}). This was corrected by removing that
guard from the altitude floor check (see Step~2 above).

%=============================================================================
\section{Change 3: parameters\_update() --- Eager Glide Airspeed Push}
%=============================================================================

\begin{lstlisting}[language=C++]
const float polar_a = _param_fw_polar_a.get();
const float polar_b = _param_fw_polar_b.get();
const float v_polar = (polar_a > FLT_EPSILON && polar_b > FLT_EPSILON)
    ? sqrtf(polar_b / polar_a) : 0.0f;
const bool polar_mode =
    (_autosoaring_control.soaring_mode == SOARING_GLIDE_POLAR);
_tecs.set_gliding_airspeed_setpoint(
    polar_mode && v_polar > FLT_EPSILON
    ? v_polar : _param_fw_glide_airspd.get());
_tecs.set_glide_i_decay(_param_fw_glide_i_decay.get());
_tecs.set_glide_stall_eas(_param_fw_airspd_stall.get());
\end{lstlisting}

\textbf{Why eager:} When Mode~2 (polar) activates mid-mission, there is
one TECS cycle between the DDS message arriving and \texttt{control\_auto()}
running. The eager push here ensures that first cycle uses the correct
polar setpoint rather than the old fixed value.

\textbf{FW\_AIRSPD\_STALL push:} \texttt{set\_glide\_stall\_eas()} stores the
stall speed so that \texttt{calcTrueAirspeedSetpoint()} can use it as the
glide airspeed floor instead of a hard-coded 3\,m/s magic constant.
\texttt{FW\_AIRSPD\_STALL} already exists in PX4 and is used by
\texttt{fw\_att\_control} and \texttt{airspeed\_selector}.
\texttt{fw\_pos\_control} now reads it and forwards it to TECS on every
parameter update.

%=============================================================================
\section{Change 4: control\_auto() --- Full Soaring Dispatch Block}
%=============================================================================

This block executes at the top of \texttt{control\_auto()} before the
inner setpoint-type switch.

\subsection{Step 1: Extract Mode Variables}

\begin{lstlisting}[language=C++]
const uint8_t soaring_cmd  = _autosoaring_control.soaring_mode;
const bool glide_cmd       = (soaring_cmd == SOARING_GLIDE_FIXED ||
                               soaring_cmd == SOARING_GLIDE_POLAR);
const bool thermal_cmd     = (soaring_cmd == SOARING_THERMAL_LOITER ||
                               soaring_cmd == SOARING_THERMAL_BANK);
const bool soaring_allowed = soaring_cmd != SOARING_OFF &&
    (_soaring_local_override || !_soaring_forbidden_latched);
\end{lstlisting}

\subsection{Step 2: Altitude Floor (FW\_ALT\_MIN)}

\begin{lstlisting}[language=C++]
// Enforced on ALL paths (DDS and CLI).
if (soaring_requested && current_altitude < alt_min) {
    if (!_soaring_forbidden_latched) {
        _soaring_forbidden_latched = true;
        if (_soaring_local_override) {
            // CLI was in control: release override so mission resumes
            _soaring_local_override  = false;
            _soaring_expect_set_mode = false;
            PX4_WARN("Soaring CLI override released: alt %.0f < FW_ALT_MIN %.0f",
                     current_altitude, alt_min);
        } else {
            PX4_WARN("Soaring disabled: alt %.0f < min %.0f",
                     current_altitude, alt_min);
        }
    }
}
if (!soaring_requested)
    _soaring_forbidden_latched = false;
\end{lstlisting}

\paragraph{Design rationale.}
The latch prevents rapid cycling near the floor altitude: once
\texttt{\_soaring\_forbidden\_latched} is set it stays set until soaring is
explicitly disabled (\texttt{soaring\_requested = false}).

\paragraph{CLI mode fix (hardware testing).}
The original code had:
\begin{lstlisting}[language=C++]
if (!_soaring_local_override && soaring_requested && ...)
\end{lstlisting}
This skipped the check when CLI was active (\texttt{\_soaring\_local\_override = true}).
The intent was ``the pilot is in control, trust them''. However, a side-effect of
the QGC-heartbeat fix (which keeps \texttt{\_soaring\_local\_override} permanently
\texttt{true} during CLI glide) was that the altitude floor was \emph{never}
triggered on hardware. The aircraft descended to ground without automatic recovery.

\textbf{Current design:} the check fires for \emph{all} paths.
When it fires during CLI override, \texttt{\_soaring\_local\_override} is also
cleared so that:
\begin{enumerate}
  \item TECS restores throttle immediately (next cycle: \texttt{effective\_glide = false}).
  \item The glide-exit block (Step~9) fires and publishes \texttt{DO\_SET\_MODE
        $\to$ AUTO\_MISSION}.
  \item The operator must re-issue \texttt{soar:glide} after regaining altitude.
\end{enumerate}

\paragraph{Altitude measurement.}
\begin{lstlisting}[language=C++]
const float current_altitude =
    _local_pos.z_global
    ? -_local_pos.z + _local_pos.ref_alt
    : _current_altitude;
\end{lstlisting}
When global-frame position is available (\texttt{z\_global = true}), the GPS
altitude is used. Otherwise the barometric estimate \texttt{\_current\_altitude}
is used. Both are above mean sea level (AMSL) — set \texttt{FW\_ALT\_MIN}
accordingly (e.g.\ home AMSL + minimum AGL).

\subsection{Step 3: Altitude Ceiling with Hysteresis}

\begin{equation}
  h \geq h_{\max} \;\Rightarrow\; \texttt{\_alt\_max\_reached} = \text{true}
\end{equation}
\begin{equation}
  h < h_{\max} - h_{\text{hyst}} \;\Rightarrow\; \texttt{\_alt\_max\_reached} = \text{false}
\end{equation}

Where $h_{\max} = $ \texttt{FW\_ALT\_MAX}, $h_{\text{hyst}} = $ \texttt{FW\_ALT\_HYST}.

When ceiling is reached during thermalling, mode degrades to glide:
\begin{equation}
  \texttt{effective\_thermal} = \text{false}, \quad \texttt{effective\_glide} = \text{true}
\end{equation}

\subsection{Step 4: Effective Mode Flags}

\begin{lstlisting}[language=C++]
const bool effective_glide =
    soaring_allowed && (glide_cmd || _alt_max_reached);
const bool effective_thermal =
    soaring_allowed && thermal_cmd && !_alt_max_reached;
const bool effective_thermal_bank =
    effective_thermal && (soaring_cmd == SOARING_THERMAL_BANK);
const bool currently_soaring = effective_glide || effective_thermal;
\end{lstlisting}

\subsection{Step 5: TECS Flag and Airspeed Injection}

\begin{lstlisting}[language=C++]
_tecs.set_gliding_mode_enabled(currently_soaring);

if (soaring_cmd == SOARING_GLIDE_FIXED
    && _autosoaring_control.airspeed_cmd > FLT_EPSILON)
    _tecs.set_gliding_airspeed_setpoint(
        _autosoaring_control.airspeed_cmd);        // Mode 1
else if (soaring_cmd == SOARING_GLIDE_POLAR) {
    const float v = sqrtf(_param_fw_polar_b.get() /
                           _param_fw_polar_a.get());
    _tecs.set_gliding_airspeed_setpoint(v);         // Mode 2
} else
    _tecs.set_gliding_airspeed_setpoint(
        _param_fw_glide_airspd.get());              // Modes 3,4

// Bank angle for load-factor stall floor
_soaring_bank_angle_cmd_deg =
    (PX4_ISFINITE(_autosoaring_control.bank_angle_cmd) &&
     _autosoaring_control.bank_angle_cmd > FLT_EPSILON)
    ? _autosoaring_control.bank_angle_cmd
    : _param_fw_thermal_bank.get();
_tecs.set_load_factor(1.f / cosf(
    radians(_soaring_bank_angle_cmd_deg)));
\end{lstlisting}

\subsection{Step 6: Thermal DO\_REPOSITION Command}

When \texttt{effective\_thermal} is true, the navigator must switch to
AUTO\_LOITER.

\subsubsection{Loiter Centre Pinning}

\begin{lstlisting}[language=C++]
const bool dds_has_pos =
    PX4_ISFINITE(_autosoaring_control.loiter_lat) &&
    fabsf(_autosoaring_control.loiter_lat) > FLT_EPSILON;
const double lat = dds_has_pos
    ? (double)_autosoaring_control.loiter_lat : curr_pos(0);
const double lon = dds_has_pos
    ? (double)_autosoaring_control.loiter_lon : curr_pos(1);
const bool first_activation = !PX4_ISFINITE(_soaring_last_lat);
const bool pos_changed  = first_activation || explicit_pos_changed;
const bool dir_changed  = !first_activation &&
    (_autosoaring_control.loiter_clockwise !=
     _soaring_last_clockwise);
\end{lstlisting}

\textbf{Bug fix:} Without pinning, DO\_REPOSITION is sent every cycle with
the drifting \texttt{curr\_pos}, causing a moving loiter centre.
The \texttt{first\_activation} flag sends exactly \emph{one}
DO\_REPOSITION. Subsequent cycles: \texttt{pos\_changed = false},
command not re-sent.

\subsubsection{T3: Pre-Roll Airspeed Safety}

On the first activation only:
\begin{equation}
  n_z = \frac{1}{\cos|\phi_{\text{cmd}}|}, \quad
  V_{\text{floor}} = 1.05 \times V_{\text{stall,level}} \times \sqrt{n_z}
\end{equation}

If $V_{\text{EAS,meas}} < V_{\text{floor}}$: delay the DO\_REPOSITION.
The aircraft stays in glide until speed is sufficient.

\subsubsection{Mission Index Latch}

Before the first DO\_REPOSITION:
\begin{lstlisting}[language=C++]
mission_s snap{};
if (_mission_sub.copy(&snap) && snap.count > 0) {
    _soaring_mission_resume_seq = (uint16_t)constrain(
        snap.current_seq, 0, snap.count - 1);
    _soaring_mission_resume_valid = true;
}
\end{lstlisting}

\subsubsection{Minimum Radius and Signed Encoding}

\begin{equation}
  R_{\min} = \frac{(V_{\text{EAS,glide}} \cdot k)^2}{g \tan|\phi_{\text{cmd}}|}
\end{equation}

\begin{equation}
  R_{\text{eff}} = \max(\texttt{loiter\_radius\_m},\; R_{\min})
\end{equation}

\begin{equation}
  \texttt{cmd.param3} = \begin{cases}
    +R_{\text{eff}} & \texttt{loiter\_clockwise} = \text{true} \\
    -R_{\text{eff}} & \texttt{loiter\_clockwise} = \text{false}
  \end{cases}
\end{equation}

\subsubsection{Complete DO\_REPOSITION Fields}

\begin{center}
\begin{tabular}{ll}
\toprule
Field & Value \\
\midrule
command & \texttt{VEHICLE\_CMD\_DO\_REPOSITION} \\
param1 & $-1$ (keep current speed) \\
param2 & 1 (switch to loiter) \\
param3 & $\pm R_{\text{eff}}$ (signed radius) \\
param4 & NaN (keep yaw) \\
param5 & loiter lat [deg] \\
param6 & loiter lon [deg] \\
param7 & current altitude [m AMSL] \\
\bottomrule
\end{tabular}
\end{center}

\subsection{Step 7: Glide DO\_SET\_MODE Heartbeat}

\begin{lstlisting}[language=C++]
if (effective_glide &&
    (now - _soaring_mode_cmd_last_us) > 1_s) {
    cmd.param2 = (float)PX4_CUSTOM_MAIN_MODE_AUTO;
    cmd.param3 = (float)PX4_CUSTOM_SUB_MODE_AUTO_MISSION;
    _pub_vehicle_command.publish(cmd);
    _soaring_mode_cmd_last_us = now;
}
\end{lstlisting}

Sent every second to keep navigator in AUTO\_MISSION while gliding.

\subsection{Step 8: Thermal Exit}

\begin{lstlisting}[language=C++]
if (_soaring_was_thermal && !effective_thermal && !effective_glide) {
    // switch back to mission
    exit_cmd.param3 = (float)PX4_CUSTOM_SUB_MODE_AUTO_MISSION;
    _pub_vehicle_command.publish(exit_cmd);
    // clear all state
    _soaring_last_lat = NAN;
    _soaring_last_lon = NAN;
    _soaring_last_clockwise = true;
    _soaring_local_override = false;
    // schedule mission resume next cycle
    if (_soaring_mission_resume_valid)
        _soaring_pending_mission_restore = true;
}
\end{lstlisting}

Mission restore is applied one cycle later, after the navigator has
confirmed it is in AUTO\_MISSION state:
\begin{lstlisting}[language=C++]
if (_soaring_pending_mission_restore
    && nav_state == NAVIGATION_STATE_AUTO_MISSION) {
    restore_cmd.command = VEHICLE_CMD_MISSION_START;
    restore_cmd.param1  = (float)_soaring_mission_resume_seq;
    _pub_vehicle_command.publish(restore_cmd);
    _soaring_pending_mission_restore = false;
}
\end{lstlisting}

%=============================================================================
\subsection{Step 9: Glide Exit --- DO\_SET\_MODE $\to$ AUTO\_MISSION}
%=============================================================================

This block is new and mirrors Step~8 for the \emph{glide} case.
When the aircraft was gliding (\texttt{\_soaring\_was\_glide = true}) and glide
just ended (altitude floor, explicit \texttt{soar:off}, or any other path),
the FMU publishes a mode-switch command so the navigator re-enters mission
flight:

\begin{lstlisting}[language=C++]
const bool soaring_was_glide_only =
    _soaring_was_glide && !_soaring_was_thermal;

if (soaring_was_glide_only && !effective_glide && !effective_thermal) {
    vehicle_command_s cmd{};
    cmd.command = VEHICLE_CMD_DO_SET_MODE;
    cmd.param1  = 1.f;                            // CUSTOM_MODE_ENABLED
    cmd.param2  = (float)PX4_CUSTOM_MAIN_MODE_AUTO;
    cmd.param3  = (float)PX4_CUSTOM_SUB_MODE_AUTO_MISSION;
    _pub_vehicle_command.publish(cmd);

    // Clear override so vehicle_command_poll() does not intercept
    // this DO_SET_MODE as a standalone "exit soaring" command.
    _soaring_local_override    = false;
    _soaring_forbidden_latched = false;
    PX4_INFO("Autosoaring: glide ended -> AUTO_MISSION");
}
\end{lstlisting}

\paragraph{Why is this needed?}
In \emph{thermal} mode the aircraft is in AUTO\_LOITER; the thermal-exit
block (Step~8) must explicitly switch back to AUTO\_MISSION.
In \emph{glide} mode the aircraft stays in AUTO\_MISSION throughout, but
the navigator has been receiving continuous \texttt{DO\_SET\_MODE(AUTO\_MISSION)}
heartbeats from the 1\,s refresh. When glide ends the navigator is already in
AUTO\_MISSION — yet republishing the command ensures the navigator
re-evaluates its setpoints and resumes waypoint tracking immediately,
rather than coasting on the last glide setpoint.

\paragraph{Guard: \texttt{soaring\_was\_glide\_only}.}
The flag \texttt{\_soaring\_was\_glide} is updated at the end of each
\texttt{control\_auto()} cycle:
\begin{lstlisting}[language=C++]
_soaring_was_glide        = effective_glide;
_soaring_was_thermal      = effective_thermal;
_soaring_was_bank_thermal = effective_thermal_bank;
\end{lstlisting}
The \texttt{!\_soaring\_was\_thermal} guard prevents a double fire when
a thermal-loiter ends and simultaneously \texttt{effective\_glide} is false:
the thermal-exit block (Step~8) already handles that case.

\paragraph{Interaction with \texttt{vehicle\_command\_poll()}.}
Clearing \texttt{\_soaring\_local\_override} before publishing ensures that
when the published \texttt{DO\_SET\_MODE} loops back through
\texttt{vehicle\_command\_poll()}, it is processed as a normal mode switch
and does not trigger the ``standalone exit soaring'' branch (which would
call \texttt{set\_gliding\_mode\_enabled(false)} redundantly).

%=============================================================================
\section{Change 5: Mode 4 --- Complete NPFG Bypass}
%=============================================================================

\subsection{Why Bypass}

If the inner switch runs \texttt{control\_auto\_loiter()} first:
\begin{itemize}
  \item NPFG computes \texttt{getAirspeedRef()} based on loiter geometry
        and wind.
  \item This modified airspeed propagates through
        \texttt{adapt\_airspeed\_setpoint()} and into TECS.
  \item TECS pitch then tracks the geometry-contaminated airspeed, not the
        polar or fixed glide airspeed.
  \item Result: airspeed oscillations and pitch-roll coupling.
\end{itemize}

\subsection{Implementation}

\begin{lstlisting}[language=C++]
if (effective_thermal_bank) {
    // Skip the switch entirely
    tecs_update_pitch_throttle(
        control_interval, current_altitude,
        _param_fw_glide_airspd.get(),   // pure glide speed
        radians(_param_fw_p_lim_min.get()),
        radians(_param_fw_p_lim_max.get()),
        _param_fw_thr_min.get(), _param_fw_thr_max.get(),
        _param_sinkrate_target.get(), _param_climbrate_target.get(),
        false);

    const float bank_sign =
        _autosoaring_control.loiter_clockwise ? 1.f : -1.f;
    const float bank_rad  =
        bank_sign * radians(fabsf(_soaring_bank_angle_cmd_deg));
    const Quatf q(Eulerf(bank_rad, get_tecs_pitch(), _yaw));
    q.copyTo(_att_sp.q_d);
    _att_sp.thrust_body[0] = get_tecs_thrust(); // = 0 (glide)

} else {
    switch (position_sp_type) { /* standard dispatch */ }
}
\end{lstlisting}

The quaternion $\mathbf{q}^*$ encodes exactly three independently sourced values:
\begin{center}
\begin{tabular}{ll}
\toprule
Component & Source \\
\midrule
$\phi^*$ & $\pm|\phi_{\text{cmd}}|$ (companion bank angle) \\
$\theta^*$ & \texttt{get\_tecs\_pitch()} (glide speed hold) \\
$\psi$ & current measured yaw (inner loop handles yaw coordination) \\
\bottomrule
\end{tabular}
\end{center}

%=============================================================================
\section{Change 6: Symmetric Throttle Ramps}
%=============================================================================

\subsection{Physical Motivation}

\paragraph{Entry (powered $\to$ glide), instant cut:}
\begin{itemize}
  \item Propwash instantly collapses, reducing wing root dynamic pressure.
  \item Nose-down pitching moment, then TECS corrects: $\pm 1.5\,\text{m/s}$
        airspeed transient.
\end{itemize}

\paragraph{Exit (glide $\to$ powered), instant restore:}
\begin{itemize}
  \item Propwash surge + pitch integrator bias: 40--50\,deg/s spike.
\end{itemize}

Both transients resolved by linear ramps of duration
$T = $ \texttt{FW\_GLIDE\_RAMP\_T} (default $3\,\text{s}$).

\subsection{Pre-Glide Thrust Cache}

Before glide starts, the current TECS thrust is saved:
\begin{lstlisting}[language=C++]
if (!currently_soaring && !_landed)
    _soaring_pre_glide_thrust = get_tecs_thrust();
\end{lstlisting}

\subsection{Entry Ramp (glide starts)}

On glide entry, \texttt{\_soaring\_entry\_ramp\_start\_us} is set.
Each cycle during entry:
\begin{equation}
  \delta_t^{\text{entry}}(t) = \delta_{t,\text{pre}} \cdot
    \max\!\left(1 - \frac{t - t_s}{T},\;0\right)
\end{equation}

After $T$ seconds: $\delta_t^{\text{entry}} = 0$.

\subsection{Exit Ramp (glide ends)}

On glide exit, \texttt{\_soaring\_ramp\_start\_us} is set.
Scale factor:
\begin{equation}
  s_{\text{exit}}(t) = \min\!\left(\frac{t - t_e}{T},\;1\right)
\end{equation}

\subsection{Final Throttle Selection}

\begin{center}
\begin{tabular}{lll}
\toprule
Condition & Thrust formula & Ramp active? \\
\midrule
Idle setpoint & $0$ & No \\
Currently soaring, entry ramp active & $\delta_{t,\text{pre}} (1 - t/T)$ & Entry \\
Currently soaring, ramp done & $0$ & None (steady glide) \\
Not soaring, exit ramp active & $\delta_{t,\text{TECS}} \cdot s_{\text{exit}}$ & Exit \\
Not soaring, normal powered & $\delta_{t,\text{TECS}}$ & None \\
Landed & $\min(\delta_{t,\text{idle}}, 1)$ & No \\
\bottomrule
\end{tabular}
\end{center}

\subsection{Physical Impact}

\begin{center}
\begin{tabular}{lll}
\toprule
Metric & Before ramp & After ramp \\
\midrule
Entry airspeed deviation & $\sim 1.5\,\text{m/s}$ & $\lesssim 0.4\,\text{m/s}$ \\
Exit pitch-rate spike & 40--50\,deg/s & $\lesssim 15\,\text{deg/s}$ \\
Throttle gradient & Step & $\leq 0.33\,\text{s}^{-1}$ (over $T = 3\,\text{s}$) \\
\bottomrule
\end{tabular}
\end{center}

%=============================================================================
\section{Complete Modification Summary}
%=============================================================================

\begin{center}
\renewcommand{\arraystretch}{1.3}
\begin{tabular}{p{0.3cm}p{3.2cm}p{4cm}p{4.2cm}}
\toprule
\# & Function / Feature & Original & Modified \\
\midrule
1 & \texttt{Run()} DDS poll & Not present & DDS poll + inhibit/watchdog \\
2 & DDS authority & N/A & CLI overrides DDS; 5\,s inhibit on \texttt{soar:off} \\
3 & Staleness watchdog & N/A & 2\,s silence $\Rightarrow$ SOARING\_OFF \\
4 & \texttt{vehicle\_command\_poll} CUSTOM\_0 & Unhandled & Encodes modes 0--4 \\
5 & DO\_SET\_MODE handling & Ignored & Pairing detection + exit logic \\
6 & \texttt{parameters\_update()} & Push FW params & Also push glide/polar airspeed \\
7 & \texttt{control\_auto()} & Mode switch only & Full soaring dispatch at top \\
8 & Altitude floor & None & Latch below \texttt{FW\_ALT\_MIN} \\
9 & Altitude ceiling & None & Hysteresis at \texttt{FW\_ALT\_MAX} \\
10 & Nav (glide) & Not commanded & DO\_SET\_MODE heartbeat 1\,s \\
11 & Nav (thermal) & Not commanded & DO\_REPOSITION once (pinned) \\
12 & Loiter direction & N/A & Signed radius, resend on flip \\
13 & Pre-roll safety & None & 5\% above banked stall \\
14 & Mission index & Not saved & Latched; restored on exit \\
15 & TECS glide flag & Never set & Set for modes 1--4 \\
16 & Airspeed setpoint & Mission speed & Mode-dependent (fixed/polar/hold) \\
17 & Mode 4 roll & None & Direct $\pm|\phi_{\text{cmd}}|$, NPFG bypassed \\
18 & Throttle entry & Instant cut & Linear ramp $T_{\text{ramp}}$ \\
19 & Throttle exit & Instant restore & Linear ramp $T_{\text{ramp}}$ \\
\bottomrule
\end{tabular}
\end{center}

\end{document}
